\begin{thebibliography}{82}
\bibitem{ref1}Misha Alekhnovich, Mark Braverman, Vitaly Feldman, Adam R. Klivans, and Toniann Pitassi. The complexity of properly learning simple concept classes. Journal of Computer and System Sciences, 74(1):16–34, 2008.
\bibitem{ref2}Eric Allender and Bireswar Das. Zero knowledge and circuit minimization. In International Symposium on Mathematical Foundations of Computer Science (MFCS), pages 25–32, 2014.
\bibitem{ref3}Eric Allender, Joshua A. Grochow, Dieter van Melkebeek, Cristopher Moore, and Andrew Morgan. Minimum circuit size, graph isomorphism, and related problems. SIAM Journal on Computing, 47(4):1339–1372, 2018.
\bibitem{ref4}Eric Allender, Lisa Hellerstein, Paul McCabe, Toniann Pitassi, and Michael E. Saks. Minimizing disjunctive normal form formulas and AC0 circuits given a truth table. SIAM Journal on Computing, 38(1):63–84, 2008.
\bibitem{ref5}Eric Allender and Shuichi Hirahara. New insights on the (non-)hardness of circuit minimization and related problems. In International Symposium on Mathematical Foundations of Computer Science (MFCS), pages 54:1–54:14, 2017.
\bibitem{ref6}Eric Allender, Dhiraj Holden, and Valentine Kabanets. The minimum oracle circuit size problem. Computational Complexity, 26(2):469–496, 2017.
\bibitem{ref7}Eric Allender, Rahul Ilango, and Neekon Vafa. The non-hardness of approximating circuit size. In International Computer Science Symposium in Russia (CSR), pages 13–24, 2019.
\bibitem{ref8}Anurag Anshu, Naresh Goud Boddu, and Dave Touchette. Quantum log-approximate-rank conjecture is also false. arXiv, 2018. arXiv:1811.10525.
\bibitem{ref9}Benny Applebaum, Boaz Barak, and David Xiao. On basing lower-bounds for learning on worst-case assumptions. In Symposium on Foundations of Computer Science (FOCS), pages 211–220, 2008.
\bibitem{ref10}Theodore P. Baker, John Gill, and Robert Solovay. Relativizations of the P =? NP question. SIAM Journal on Computing, 4(4):431–442, 1975.
\bibitem{ref11}Avrim Blum, Merrick L. Furst, Michael J., and Richard J. Lipton. Cryptographic primitives based on hard learning problems. In International Cryptology Conference (CRYPTO), pages 278–291, 1993.
\bibitem{ref12}Avrim L. Blum and Ronald L. Rivest. Training a 3-node neural network is np-complete. Neural Networks, 5(1):117–127, 1992.
\bibitem{ref13}Anselm Blumer, Andrzej Ehrenfeucht, David Haussler, and Manfred K. Warmuth. Occam’s razor. Information Processing Letters, 24(6):377–380, 1987.
\bibitem{ref14}Andrej Bogdanov and Luca Trevisan. Average-case complexity. Foundations and Trends in Theoretical Computer Science, 2(1), 2006.
\bibitem{ref15}Dan Boneh and Richard J. Lipton. Amplification of weak learning under the uniform distribu- tion. In Conference on Learning Theory (COLT), pages 347–351, 1993.
\bibitem{ref16}Joan Boyar, Philip Matthews, and René Peralta. On the shortest linear straight-line program for computing linear forms. In International Symposium on Mathematical Foundations of Computer Science (MFCS), pages 168–179, 2008.
\bibitem{ref17}Marco L. Carmosino, Russell Impagliazzo, Valentine Kabanets, and Antonina Kolokolova. Learning algorithms from natural proofs. In Conference on Computational Complexity (CCC), pages 10:1–10:24, 2016.
\bibitem{ref18}Amit Chakrabarti and Oded Regev. An optimal lower bound on the communication complexity of gap-Hamming-distance. SIAM Journal on Computing, 41(5):1299–1317, 2012. doi:10. 1137/120861072.
\bibitem{ref19}Arkadev Chattopadhyay, Nikhil S Mande, and Suhail Sherif. The log-approximate-rank conjecture is false. In Symposium on Theory of Computing (STOC), pages 42–53. ACM, 2019.
\bibitem{ref20}Mahdi Cheraghchi, Valentine Kabanets, Zhenjian Lu, and Dimitrios Myrisiotis. Circuit lower bounds for MCSP from local pseudorandom generators. Electronic Colloquium on Computational Complexity (ECCC), 26:22, 2019.
\bibitem{ref21}Gil Cohen and Igor Shinkar. The complexity of DNF of parities. In Innovations in Theoretical Computer Science (ITCS), pages 47–58, 2016.
\bibitem{ref22}Sebastian Czort. The complexity of minimizing disjunctive normal form formulas. Master’s Thesis, University of Aarhus, 1999.
\bibitem{ref23}Uriel Feige. A threshold of ln n for approximating set cover. Journal of the ACM, 45(4):634–652, 1998.
\bibitem{ref24}Uriel Feige and Joe Kilian. Zero knowledge and the chromatic number. Journal of Computer and System Sciences, 57(2):187–199, 1998.
\bibitem{ref25}Magnus Gausdal Find, Alexander Golovnev, Edward A. Hirsch, and Alexander S. Kulikov. A better-than-3n lower bound for the circuit complexity of an explicit function. In Symposium on Foundations of Computer Science (FOCS), pages 89–98, 2016.
\bibitem{ref26}Michael R. Garey and David S. Johnson. Computers and Intractability: A Guide to the Theory of NP-Completeness. W. H. Freeman, 1979.
\bibitem{ref27}Alexander Golovnev, Rahul Ilango, Russell Impagliazzo, Valentine Kabanets, Antonina Koloko- lova, and Avishay Tal. AC0 [p] lower bounds against MCSP via the coin problem. Electronic Colloquium on Computational Complexity (ECCC), 26:18, 2019.
\bibitem{ref28}Mika Göös and Toniann Pitassi. Communication lower bounds via critical block sensitivity. SIAM Journal on Computing, 47(5):1778–1806, 2018. doi:10.1137/16M1082007.
\bibitem{ref29}Gary D. Hachtel and Fabio Somenzi. Logic synthesis and verification algorithms. Springer, 2006.
\bibitem{ref30}Thomas R. Hancock, Tao Jiang, Ming Li, and John Tromp. Lower bounds on learning decision lists and trees. Information and Computation, 126(2):114–122, 1996.
\bibitem{ref31}Johan Hastad. Clique is hard to approximate within $n^{1-\varepsilon}$ . In Symposium on Foundations of Computer Science (FOCS), pages 627–636, 1996.
\bibitem{ref32}Shuichi Hirahara. Non-black-box worst-case to average-case reductions within NP. In Sym- posium on Foundations of Computer Science (FOCS), pages 247–258, 2018.
\bibitem{ref33}Shuichi Hirahara, Igor Carboni Oliveira, and Rahul Santhanam. NP-hardness of minimum circuit size problem for OR-AND-MOD circuits. In Computational Complexity Conference (CCC), pages 5:1–5:31, 2018.
\bibitem{ref34}Shuichi Hirahara and Rahul Santhanam. On the average-case complexity of MCSP and its variants. In Computational Complexity Conference (CCC), pages 7:1–7:20, 2017.
\bibitem{ref35}Shuichi Hirahara and Osamu Watanabe. Limits of minimum circuit size problem as oracle. In Conference on Computational Complexity (CCC), pages 18:1–18:20, 2016.
\bibitem{ref36}John M. Hitchcock and Aduri Pavan. On the NP-completeness of the minimum circuit size problem. In Conference on Foundation of Software Technology and Theoretical Computer Science (FSTTCS), pages 236–245, 2015.
\bibitem{ref37}Sangxia Huang. Improved hardness of approximating chromatic number. In International Con- ference on Approximation Algorithms for Combinatorial Optimization Problems (APPROX), pages 233–243, 2013.
\bibitem{ref38}Rahul Ilango. AC0 [p] lower bounds and NP-hardness for variants of MCSP. Electronic Colloquium on Computational Complexity (ECCC), 26:21, 2019.
\bibitem{ref39}Russell Impagliazzo, Valentine Kabanets, and Ilya Volkovich. The Power of Natural Properties as Oracles. In Computational Complexity Conference (CCC), volume 102, pages 7:1–7:20, 2018.
\bibitem{ref40}Kazuo Iwama and Hiroki Morizumi. An explicit lower bound of 5n − o(n) for boolean circuits. In International Symposium on Mathematical Foundations of Computer Science (MFCS), pages 353–364, 2002.
\bibitem{ref41}J. Stephen Judd. Learning in networks is hard. In International Conference on Neural Networks (ICNN), volume 2, pages 685–692, 1987.
\bibitem{ref42}Stasys Jukna. Boolean function complexity: advances and frontiers, volume 27. Springer, 2012.
\bibitem{ref43}Valentine Kabanets and Jin-yi Cai. Circuit minimization problem. In Symposium on Theory of Computing (STOC), pages 73–79, 2000.
\bibitem{ref44}Maurice Karnaugh. The map method for synthesis of combinational logic circuits. Transactions of the American Institute of Electrical Engineers, Part I: Communication and Electronics, 72(5):593–599, 1953.
\bibitem{ref45}Michael J. Kearns and Leslie G. Valiant. Cryptographic limitations on learning boolean formulae and finite automata. J. ACM, 41(1):67–95, 1994. doi:10.1145/174644.174647.
\bibitem{ref46}Michael J. Kearns and Umesh V. Vazirani. An Introduction to Computational Learning Theory. MIT Press, 1994.
\bibitem{ref47}Subhash Khot. Improved inapproximability results for maxclique, chromatic number and approximate graph coloring. In Symposium on Foundations of Computer Science (FOCS), pages 600–609, 2001.
\bibitem{ref48}Subhash Khot and Rishi Saket. Hardness of minimizing and learning DNF expressions. In Symposium on Foundations of Computer Science (FOCS), pages 231–240, 2008.
\bibitem{ref49}Ker-I Ko. On the complexity of learning minimum time-bounded Turing machines. SIAM Journal on Computing, 20(5):962–986, 1990.
\bibitem{ref50}Pravesh K. Kothari and Roi Livni. Improper learning by refuting. In Innovations in Theoretical Computer Science (ITCS), pages 55:1–55:10, 2018.
\bibitem{ref51}Jan Krajícek. Forcing with Random Variables and Proof Complexity. Cambridge University Press, 2011.
\bibitem{ref52}Eyal Kushilevitz. Communication complexity. In Advances in Computers, volume 44, pages 331–360. Elsevier, 1997.
\bibitem{ref53}Eyal Kushilevitz and Noam Nisan. Communication Complexity. Cambridge University Press, 1997.
\bibitem{ref54}Eyal Kushilevitz and Enav Weinreb. On the complexity of communication complexity. In Symposium on Theory of Computing (STOC), pages 465–474, 2009.
\bibitem{ref55}Leonid Levin. Universal sequential search problems. Problemy Peredachi Informatsii, 9(3):115– 116, 1973.
\bibitem{ref56}László Lovász and Michael Saks. Lattices, mobius functions and communications complexity. In Symposium on Foundations of Computer Science (FOCS), pages 81–90, 1988.
\bibitem{ref57}Carsten Lund and Mihalis Yannakakis. On the hardness of approximating minimization problems. Journal of the ACM, 41(5):960–981, 1994.
\bibitem{ref58}William J. Masek. Some NP-complete set covering problems. Unpublished Manuscript, 1979.
\bibitem{ref59}Edward J. McCluskey. Introduction to the theory of switching circuits. McGraw-Hill, 1965.
\bibitem{ref60}Dylan M. McKay, Cody D. Murray, and R. Ryan Williams. Weak lower bounds on resource- bounded compression imply strong separations of complexity classes. In Symposium on Theory of Computing (STOC), 2019.
\bibitem{ref61}Moritz Müller and Ján Pich. Feasibly constructive proofs of succinct weak circuit lower bounds. Electronic Colloquium on Computational Complexity (ECCC), 24:144, 2017.
\bibitem{ref62}Cody D. Murray and Richard Ryan Williams. On the (non) NP-hardness of computing circuit complexity. In Conference on Computational Complexity (CCC), pages 365–380, 2015.
\bibitem{ref63}Igor Carboni Oliveira, Ján Pich, and Rahul Santhanam. Hardness magnification near state-of- the-art lower bounds. In Conference on Computational Complexity (CCC), 2019.
\bibitem{ref64}Igor Carboni Oliveira and Rahul Santhanam. Conspiracies between learning algorithms, circuit lower bounds, and pseudorandomness. In Computational Complexity Conference (CCC), pages 18:1–18:49, 2017.
\bibitem{ref65}Igor Carboni Oliveira and Rahul Santhanam. Hardness magnification for natural problems. In Symposium on Foundations of Computer Science (FOCS), pages 65–76, 2018.
\bibitem{ref66}James Orlin. Contentment in graph theory: covering graphs with cliques. Indagationes Mathematicae, 80(5):406–424, 1977.
\bibitem{ref67}Leonard Pitt and Leslie G. Valiant. Computational limitations on learning from examples. Journal of the ACM, 35(4):965–984, 1988.
\bibitem{ref68}Pavel Pudlák, Vojtech Rödl, and Petr Savický. Graph complexity. Acta Informatica, 25(5):515– 535, 1988.
\bibitem{ref69}Netanel Raviv.     Truth table minimization of computational models.           CoRR/arXiv, abs/1306.3766, 2013. arXiv:1306.3766.
\bibitem{ref70}Alexander A. Razborov and Steven Rudich. Natural proofs. Journal of Computer and System Sciences, 55(1):24–35, 1997.
\bibitem{ref71}Ronald L. Rivest. Learning decision lists. Machine Learning, 2(3):229–246, 1987.
\bibitem{ref72}J. Paul Roth and Richard M. Karp. Minimization over boolean graphs. IBM Journal of Research and Development, 6(2):227–238, 1962.
\bibitem{ref73}Christoph Scholl. Functional decomposition with applications to FPGA synthesis. Kluwer Academic Publishers, 2001.
\bibitem{ref74}Makrand Sinha and Ronald de Wolf. Exponential separation between quantum communication and logarithm of approximate rank. arXiv, 2018. arXiv:1811.10090.
\bibitem{ref75}Boris A Trakhtenbrot. A survey of Russian approaches to perebor (brute-force search) algorithms. Annals of the History of Computing, 6(4):384–400, 1984.
\bibitem{ref76}Luca Trevisan. Non-approximability results for optimization problems on bounded degree instances. In Symposium on Theory of Computing (STOC), pages 453–461, 2001.
\bibitem{ref77}Christopher Umans, Tiziano Villa, and Alberto L. Sangiovanni-Vincentelli. Complexity of two-level logic minimization. IEEE Trans. on CAD of Integrated Circuits and Systems, 25(7):1230–1246, 2006.
\bibitem{ref78}Salil P. Vadhan. On learning vs. refutation. In Conference on Learning Theory (COLT), pages 1835–1848, 2017.
\bibitem{ref79}Marcin Wrochna and Stanislav Živnỳ. Improved hardness for H-colourings of G-colourable graphs. arXiv, 2019. arXiv:1907.00872.
\bibitem{ref80}Andrew Chi-Chih Yao. Some complexity questions related to distributive computing (pre- liminary report). In Symposium on Theory of Computing (STOC), pages 209–213, 1979. doi:10.1145/800135.804414.
\bibitem{ref81}Andrew Chi-Chih Yao. Theory and applications of trapdoor functions (extended abstract). In Symposium on Foundations of Computer Science (FOCS), pages 80–91, 1982.
\bibitem{ref82}David Zuckerman. Linear degree extractors and the inapproximability of max clique and chromatic number. In Symposium on Theory of Computing (STOC), pages 681–690, 2006.
\end{thebibliography}
